\chapter*{ЗАКЛЮЧЕНИЕ}
\addcontentsline{toc}{chapter}{ЗАКЛЮЧЕНИЕ}

В работе предложены и исследованы методы оптимизации архитектуры IFRNet, направленные на повышение вычислительной эффективности при сохранении качества интерполяции видеокадров.

\subsection*{Основные результаты}

\begin{enumerate}
	\item Предложенные методы (сокращение числа каналов и переход к YUV420) обеспечили снижение вычислительной сложности в 3,8--3,9 раза по FLOPs и уменьшение числа параметров в 4,96 раза (с 4,96~M до 1,00~M); ускорение инференса составило более чем в два раза (с 0,032 до 0,014--0,015~с/кадр).

	\item Наилучший баланс между эффективностью и качеством показала модель \texttt{IFRNet\_RGB2+LPIPS}: при снижении FLOPs в 3,8 раза она обеспечивает наилучшее среди всех моделей перцептуальное качество (LPIPS = 0,0361, DISTS = 0,0222), превосходящее в том числе базовую архитектуру (LPIPS = 0,0503, DISTS = 0,0342), при близких пиксельных метриках (PSNR = 35,129~дБ, SSIM = 0,9774).

	\item Добавление перцептуальной функции потерь LPIPS компенсирует снижение визуального качества, вызванное архитектурной оптимизацией: относительно модели \texttt{IFRNet\_RGB2} без LPIPS улучшение составило около 31~\% по метрике LPIPS и около 39~\% по метрике DISTS.

	\item Переход к цветовому пространству YUV420 обеспечивает минимальные значения FLOPs и времени инференса, однако уступает варианту RGB2 по перцептуальным метрикам; его применение целесообразно при интеграции в конвейеры обработки видео, изначально использующие представление YUV.
\end{enumerate}

\subsection*{Перспективы дальнейших исследований}

Дальнейшие исследования могут быть направлены на:
\begin{itemize}
	\item разработку гибридной стратегии, адаптивно выбирающей между режимами RGB2 и YUV420 в зависимости от характеристик сцены (уровня детализации текстур, интенсивности движения), что позволит динамически балансировать вычислительную эффективность и качество интерполяции;
	\item применение других перцептуальных функций потерь (например, DISTS) в процессе обучения для дальнейшего повышения визуального качества оптимизированных моделей;
	\item адаптацию предложенных моделей для мобильных и встраиваемых устройств с использованием аппаратного ускорения (TensorRT, TFLite) и методов квантования.
\end{itemize}
